\documentclass{amsart}

\usepackage{keytheorems}
\usepackage{amsthm}
\usepackage{color}

\newkeytheorem{problem}
\newtheorem*{proposition}{Proposition}
\newcommand{\todo}{\textcolor{red}{TO-DO.} }
\newcommand{\temp}{\textcolor{red}{TEMPORARY.} }

\newcommand{\set}[1]{\left\{#1\right\}}
\newcommand{\powset}{{\cal P}}
\newcommand{\img}{_{*}}
\newcommand{\pimg}{^{*}}
\newcommand{\abs}[1]{\left|#1\right|}
\newcommand{\ceil}[1]{\lceil #1 \rceil}
\newcommand{\floor}[1]{\lfloor #1 \rfloor}
\newcommand{\half}{\frac{1}{2}}

\newcommand{\R}{\mathbb{R}}
\newcommand{\Z}{\mathbb{Z}}
\newcommand{\C}{\mathbb{C}}
\newcommand{\N}{\mathbb{N}}
\newcommand{\Q}{\mathbb{Q}}

\numberwithin{equation}{section}

\begin{document}

	\title[Math 251W: Homework \#13]{
		Math 251W: Foundations of Advanced Mathematics \\
		\textmd{Homework \#13: Portfolio Problems \S 6.5 \& 6.6}}
	\author{Carmen Beltran, Brendan Butler, Teddy Wachtler}
	\date{\today}

	\maketitle

	\begin{problem}[manual-num=65]
		\temp
		\begin{proposition}
			Prove that for all sets $A, B$, and $X$ so that $X \subseteq A$ and functions $f : A \rightarrow B$, $f$ being injective implies $|X| = |f \img (X)|$.
		\end{proposition}

		\begin{proof}
			Let $A, B$, and $X$ be arbitrary sets so that $X \subseteq A$.
			Let $f : A \rightarrow B$ be an arbitrary injective function.
			Let $g : X \rightarrow f \img (X)$ be the function defined by $g(x) = f(x)$.

			Let $x_0$ be an arbitrary element of $X$.
			By substitution, we see that $g(x_0) = f(x_0)$.
			By the definition of $f \img (X)$, since $x_0 \in X$, we see that $f(x_0) \in f \img (X)$.
			By substitution, we see that $g(x_0) \in f \img (X)$.
			By the definition of $g$, we see that $g$ is well-defined.

			Let $x_1$ and $x_2$ be arbitrary elements of $X$ so that $g(x_1) = g(x_2)$.
			By substitution, we see that $f(x_1) = f(x_2)$.
			By the definition of injective functions, since $f$ is injective, we see that $x_1 = x_2$.
			By the definition of injective functions, we see that $g$ is injective.

			Let $y$ be an arbitrary element in $f \img (X)$.
			By the definition of $f \img (X)$, there exists a particular element $x_3 \in X$ so that $f(x_3) = y$.
			By substitution, since $g(x) = f(x)$, we see that $g(x_3) = y$.
			By the definition of surjective functions, we see that $g$ is surjective.

			By the definition of bijective functions, since $g$ is injective and surjective, we see that $g$ is bijective.
			Since there exists a bijective function between $X$ and $f \img (X)$, we see that $\abs{X} = \abs{f \img (X)}$.
			Therefore, we see that for all sets $A, B$, and $X$ so that $X \subseteq A$ and functions $f : A \rightarrow B$, $f$ being injective implies $|X| = |f \img (X)|$.
		\end{proof}
	\end{problem}

	\begin{problem}[manual-num=64]
		\begin{proposition}
			Prove that the set of all integers that are a multiple of 5 has the same cardinality as the set of all integers.
		\end{proposition}

		\begin{proof}
			Let $X$ be the particular set defined by $X = \set{5x : x \in \Z}$.
			Let $f : \Z \rightarrow X$ be the function defined by $f(x) = 5x$.

			Let $x_0$ be an arbitrary element of $\Z$.
			By substitution, we see that $f(x_0) = 5x_0$.
			By the definition of $X$, we see that $f(x_0) \in X$.
			By the definition of $f$, we see that $f$ is well-defined.

			Let $x_1$ and $x_2$ be arbitrary elements of $\Z$ so that $f(x_1) = f(x_2)$.
			By substitution, we see that $5x_1 = 5x_2$.
			By division and substitution, we see that $x_1 = x_2$.
			By the definition of injective functions, we see that $f$ is injective.

			Let $y$ be an arbitrary element in $X$.
			By the definition of $X$, we see that there exists $x_3 \in \Z$ so that $y = 5x_3$.
			By substitution, we see that $f(x_3) = 5x_3 = y$.
			By the definition of surjective functions, we see that $f$ is surjective.

			By the definition of bijective functions, since $f$ is injective and surjective, we see that $f$ is bijective.
			Since there exists a bijective function between $\Z$ and $X$, we see that $\abs{\Z} = \abs{X}$.
			Therefore, for the set $X = \set{5x : x \in \Z}$, we see that $\abs{\Z} = \abs{X}$.
		\end{proof}
	\end{problem}

	\begin{problem}[manual-num=66]
		\begin{proposition}
			Let $a, b, c, d \in \R$.
			Suppose $a < b$ and $c < d$.
			Use the Schroeder-Bernstein Theorem to prove that $\abs{[a, b)} = \abs{[c, d]}$.
		\end{proposition}

		\begin{proof}
			Let $a, b, c$, and $d$ be arbitrary real numbers.
			Suppose $a < b$ and $c < d$.

			Let $f : [a, b) \rightarrow [c, d]$ be the function defined by $f(x) = \frac{(x - a)(d - c)}{(b - a)} + c$.
			Since $f$ is strictly increasing, we see that $f(a) \leq f(x) < f(b)$ for all $x \in [a, b)$.
			By algebra, for all $x \in [a, b)$, we see that:
			\begin{align*}
				f(a) &\leq f(x) < f(b) \\
				\frac{0 \cdot (d - c)}{(b - a)} + c &\leq f(x) < (d - c) + c \\
				c &\leq f(x) < d \\
				c &\leq f(x) \leq d.
			\end{align*}
			Since $f(x) \in [c, d]$ for all $x \in [a, b)$, we see that $f$ is well-defined.
			Let $x_1$ and $x_2$ be arbitrary elements of $[a, b)$ so that $f(x_1) = f(x_2)$.
			By algebra, we see that:
			\begin{align*}
				f(x_1) &= f(x_2) \\
				\frac{(x_1 - a)(d - c)}{(b - a)} + c &= \frac{(x_2 - a)(d - c)}{(b - a)} + c \\
				(x_1 - a)(d - c) &= (x_2 - a)(d - c) \\
				(x_1 - a) &= (x_2 - a) \quad \text{(since $c < d$ implies $d - c \neq 0)$} \\
				x_1 &= x_2.
			\end{align*}
			By the definition of injective functions, we see that $f$ is injective.

			Let $g : [c, d] \rightarrow [a, b)$ be the function defined by $g(y) = \frac{(y - c)(b - a)}{2 \cdot (d - c)} + a$.
			Since $g$ is strictly increasing, we see that $g(c) \leq g(y) \leq g(d)$ for all $y \in [c, d]$.
			By algebra, for all $y \in [c, d]$, we see that:
			\begin{align*}
				g(c) &\leq g(y) \leq g(d) \\
				\frac{0 \cdot (b - a)}{2(d - c)} + a &\leq g(y) \leq \frac{(b - a)}{2} + a \\
				a &\leq g(y) \leq \frac{a + b}{2}. \\
				a &\leq g(y) < b \quad \text{(since $a < b$ implies $\frac{a + b}{2} < b$)}.
			\end{align*}
			Since $g(y) \in [a, b)$ for all $y \in [c, d]$, we see that $g$ is well-defined.
			Let $y_1$ and $y_2$ be arbitrary elements of $[c, d]$ so that $g(y_1) = g(y_2)$.
			By algebra, we see that:
			\begin{align*}
				g(y_1) &= g(y_2) \\
				\frac{(y_1 - c)(b - a)}{2(d - c)} + a &= \frac{(y_2 - c)(b - a)}{2(d - c)} + a \\
				(y_1 - c)(b - a) &= (y_2 - c)(b - a) \\
				(y_1 - c) &= (y_2 - c) \quad \text{(since $a < b$ implies $b - a \neq 0$)} \\
				y_1 &= y_2.
			\end{align*}
			By the definition of injective functions, we see that $g$ is injective.

			By the Schroeder-Bernstein Theorem, since there exists injective functions $f : [a, b) \rightarrow [c, d]$ and $g : [c, d] \rightarrow [a, b)$, we see that $\abs{[a, b)} = \abs{[c, d]}$.
			Therefore, for all real numbers $a, b, c$, and $d$, we see that $a < b$ and $c < d$ implies $\abs{[a, b)} = \abs{[c, d]}$.
		\end{proof}
	\end{problem}

	\begin{problem}[manual-num=67]
		\begin{proposition}
			Let $X$ be an arbitrary countably infinite set.
			Prove that there exists a function $f : X \rightarrow X$  that is injective but not surjective.
		\end{proposition}

		\begin{proof}
			Let $X$ be an arbitrary countably infinite set.
			By the definition of countably infinite sets, we see that $X = \set{x_n : n \in \N^+}$ and that there exists a bijective function $g : \N^+ \rightarrow X$ defined by $g(n) = x_n$.
			Let $f : X \rightarrow X$ be the function defined by $f(x_n) = g(n + 1)$.

			Let $a$ be an arbitrary positive natural number.
			By the definition of $X$, we see that $x_a$ is an arbitrary element of $X$.
			By substitution, we see that $f(x_a) = g(a + 1)$.
			Since $g$ is bijective, we see that $g(a + 1)$ is a unique element in $X$.
			By the definition of $f$, we see that $f$ is well-defined.

			Let $b$ and $c$ be arbitrary positive natural numbers.
			By the definition of $X$, we see that $x_b$ and $x_c$ are arbitrary elements of $X$.
			Suppose $f(x_b) = f(x_c)$.
			By substitution, we see that $g(b + 1) = g(c + 1)$.
			Since $g$ is bijective, we see that $b + 1 = c + 1$.
			By subtraction, we see that $b = c$.
			By substitution, we see that $x_b = x_c$.
			By the definition of injective functions, we see that $f$ is injective.

			By way of contradiction, suppose there exists a positive natural number $d$ so that $f(x_d) = x_1$.
			By substitution, we see that $g(d + 1) = g(1)$.
			Since $g$ is bijective, we see that $d + 1 = 1$.
			By subtraction, we see that $d = 0$.
			Since $0$ is not a positive natural number, we see that $f(x_n) \neq x_1$ for all positive natural numbers $n$ and elements $x_n \in X$.
			By the definition of surjective functions, we see that $f$ is not surjective.

			Therefore, we see that for all countably infinite sets $X$, there exists a function $f : X \rightarrow X$ that is injective but not surjective.
		\end{proof}
	\end{problem}

\end{document}